The project is strategically aligned with the long-term development goals of the consortium partners, addressing specific institutional and market needs:

\begin{itemize}
  \item \textbf{Springer Nature:} 
  The project integrates directly into their transition from a traditional publisher to a comprehensive \textbf{Open Science platform provider}. By treating high-fidelity 3D models and XR assets as citable research outputs, Springer Nature addresses the market need for standardized digital heritage documentation and open access revenue models.
  \item \textbf{Arts et Métiers (CAPLAB / Laval Virtual):}
  A core strategic goal is to expand their \textbf{4D reality capture technologies} from industrial and healthcare sectors into the cultural heritage market. The project addresses their need for technological validation (increasing TRL) in complex, non-controlled outdoor environments, which is essential for their market segment expansion.
  \item \textbf{University Partners (ELTE, Sapienza, NKUA, Sorbonne):}
  The project supports their internationalization and \textbf{Digital Humanities development strategies}. It addresses the need for interdisciplinary research frameworks that combine classical archaeology with AI-based processing and climate risk assessments.
\end{itemize}

This synergy ensures that project outcomes are not only scientifically relevant but are also embedded in the sustainable business and research models of all partners beyond the funding cycle.